% \begin{theorem}\label{old-QESAT_poly_LPCP}\lean{QESAT_poly_LPCP}\uses{QESAT, LPCP}\leanok
% \[
%   \mathrm{QESAT}(\F_2,n) \in
%   \mathrm{LPCP}[
%     \varepsilon_c = 0,
%     \varepsilon_s = 3/4,
%     \Sigma = \F_2,
%     \ell = n + n^2,
%     q = 4,
%     r = |x| + 2*n
%   ].
% \]
% \end{theorem}
% \begin{proof}%\uses{} tensor stuff

% \noindent Let $(p_1, \dots, p_m)$ be an instance of $\text{QESAT}(\mathbb{F})$ with $n$ variables. \\
% The LPCP verifier expects a proof $\pi = (a, b) \in \mathbb{F}^{n+n^2}$ and works as follows:

% \vspace{1em}
% \noindent \textbf{Verifier} $V^{(a,b)}(p_1, \dots, p_m)$:
% \begin{enumerate}
%     \item Sample $r \gets \mathbb{F}^m$ and $s, t \gets \mathbb{F}^n$.
%     \item Define 
%     \[
%     M := \begin{bmatrix} \text{coeff}(p_1) \\ \vdots \\ \text{coeff}(p_m) \end{bmatrix} \in \mathbb{F}^{m \times (n+n^2)} 
%     \quad \text{and} \quad 
%     c := \begin{bmatrix} -\text{const}(p_1) \\ \vdots \\ -\text{const}(p_m) \end{bmatrix} \in \mathbb{F}^m.
%     \]
%     \textit{(Note: $\text{coeff}(p_i) :=$ non-constant coeffs of $p_i$, and $\text{const}(p_i) :=$ constant coeff of $p_i$)}
    
%     \item \textit{(linear equation test)} $\rightarrow$ Query $(a,b)$ at $M^T r$ and check that $\langle a \parallel b, M^T r \rangle = \langle c, r \rangle$.
    
%     \item \textit{(tensor test)} $\rightarrow$ Query $b$ at $\text{flat}(s \otimes t)$, $a$ at $s$ and $t$, and check that 
%     \[
%     \langle b, \text{flat}(s \otimes t) \rangle = \langle a, s \rangle \cdot \langle a, t \rangle.
%     \]
% \end{enumerate}

% \vspace{1em}
% \noindent \textbf{Completeness:} 
% Suppose $p_1(a) = \dots = p_m(a) = 0$ and set $b := \text{flat}(a \otimes a)$. \\
% Then 
% \[
% \Pr_{s,t} \left[ \langle b, \text{flat}(s \otimes t) \rangle = \langle a, s \rangle \cdot \langle a, t \rangle \right] = 1
% \]
% and 
% \[
% \Pr_r \left[ \langle a \parallel b, M^T r \rangle = \langle c, r \rangle \right] = 1 \quad \left(\text{since } M \begin{bmatrix} a \\ \text{flat}(a \otimes a) \end{bmatrix} = c \right).
% \]

% \vspace{1em}
% \noindent \textbf{Soundness:}
% Suppose $\forall a \in \mathbb{F}^n \ \exists i \in [m] \text{ s.t. } p_i(a) \neq 0$. Fix any $\pi = (a, b)$. \\
% Either:
% \begin{enumerate}
%     \item[(i)] $b \neq \text{flat}(a \otimes a) \implies$ tensor test passes w.p. $\leq \frac{2|\mathbb{F}|-1}{|\mathbb{F}|^2}$
%     \item[OR (ii)] $b = \text{flat}(a \otimes a)$ and $M \begin{bmatrix} a \\ b \end{bmatrix} \neq c \implies$ linear equation test passes w.p. $\leq \frac{1}{|\mathbb{F}|}$
% \end{enumerate}

% \end{proof}


\begin{definition}
\label{QESAT.linearization}
\uses{QESAT}
Let \(x=(p_1,\dots,p_m)\) be a \(\QESAT(\F_2,n)\) instance. For each quadratic
polynomial \(p_i\), form a row \(M_i\in\F_2^{n+n^2}\) by replacing every
linear monomial \(X_j\) by the variable \(a_j\) and every quadratic monomial
\(X_jX_k\) by the tensor variable \(b_{jk}\). Let
    \[
    M_{x} := \begin{bmatrix} M_1 \\ \vdots \\ M_m \end{bmatrix} \in \F_2^{m \times (n+n^2)} 
    \quad \text{and} \quad 
    c_{x} := \begin{bmatrix} -\operatorname{const}(p_1) \\ \vdots \\ -\operatorname{const}(p_m) \end{bmatrix} \in \F_2^m,
    \]
where \(\operatorname{const}(p_i)\) is the constant term of \(p_i\).
Then \(M_{x}(a\parallel b)=c_{x}\) expresses the original equations after
replacing \(a_ja_k\) by \(b_{jk}\).
\end{definition}

\begin{definition}
\label{QESAT.verifier}
\lean{QESAT.verifier}
\uses{LPCPVerifier,QESAT.linearization,LINEQ.verifier,TENSORQ.selfVerifier}
\leanok
Let \(x=(p_1,\dots,p_m) \in R_{2}^{m}\) be a QESAT instance as defined in \ref{QESAT}. Let \(V_{\QESAT(\F_2,n)}^{\pi}(x)\) be the LPCP verifier 
with oracle access to \(\pi=a\parallel b\in\F_2^{n+n^2}\), defined as follows:
\begin{enumerate}
  \item Sample $r \gets \F_{2}^m$ and $s, t \gets \F_{2}^n$.
  \item Construct \((M_{x},c_{x}) \in \F_2^{m \times (n+n^2)} \times \F_2^m\) as in \ref{QESAT.linearization}.
  \item Check that \(V_{\LINEQ(\F_2,m,n+n^2)}^{\langle \pi, \cdot \rangle } (M_{x},c_{x}) = 1\).
  \item Query \(\pi\) at
  \((s,0)\), \((t,0)\), and \((0,\mathrm{flat}(s\otimes t))\) and check that
  \[
    \langle b,\mathrm{flat}(s\otimes t)\rangle
    =\langle a,s\rangle\langle a,t\rangle .
  \]
  Equivalently, this last step is the tensor self-check of \ref{TENSORQ.selfVerifier} specialized to \(\F_2\).
\end{enumerate}
\end{definition}

\begin{lemma}[Completeness]
\label{QESAT.verifier_completeness}
\lean{QESAT.verifier_correct}
\uses{QESAT.verifier,QESAT.linearization,LINEQ.verifier,TENSORQ.tensor_check_complete}
\leanok
If \(x\in \QESAT(\F_2,n)\), then exists a proof
\(\pi\in\F_2^{n+n^2}\) such that \[\Pr\!\left[V_{\QESAT(\F_2,n)}^{\pi}(x)=1\right]=1.\]
\end{lemma}
\begin{proof}
Let \(a\in\F_2^n\) satisfy all equations and set
\(b=\mathrm{flat}(a\otimes a)\). By construction of
\ref{QESAT.linearization}, \(M_x(a\parallel b)=c_x\), so the LINEQ check
accepts for every \(r\). The tensor self-check accepts for every \(s,t\) by
\ref{TENSORQ.tensor_check_complete}.
\end{proof}

\begin{lemma}[Soundness]
\label{QESAT.verifier_soundness}
\lean{QESAT.verifier_correct}
\uses{QESAT.verifier,QESAT.linearization,LINEQ.linear_form_uniform_prob_mul_card_le_one,TENSORQ.selfVerifier_soundness_after_sampling}
\leanok
If \(x\notin \QESAT(\F_2,n)\), then for every proof
\(\tilde\pi\in\F_2^{n+n^2}\),
\[\Pr\!\left[V_{\QESAT(\F_2,n)}^{\tilde\pi}(x)=1\right]\leq 3 / 4.\]
\end{lemma}
\begin{proof}
If some \(p_i\) has total degree \(>2\), the verifier rejects immediately, so
the acceptance probability is \(0\). Assume now that every \(p_i\) is quadratic
and write \(\pi=a\parallel b\). If \(b\neq\mathrm{flat}(a\otimes a)\), then
the tensor self-check accepts with probability at most \(3/4\) by
\ref{TENSORQ.selfVerifier_soundness_after_sampling} specialized to \(\F_2\), and
the full verifier accepts only if that check accepts.

Otherwise \(b=\mathrm{flat}(a\otimes a)\). Since \(x\notin\QESAT(\F_2,n)\),
the assignment \(a\) does not satisfy all equations. By
\ref{QESAT.linearization}, \(M_x(a\parallel b)\neq c_x\). The LINEQ check
therefore accepts with probability at most \(1/|\F_2|=1/2\) by
\ref{LINEQ.linear_form_uniform_prob_mul_card_le_one}. Again the full verifier
accepts only if this check accepts, so the acceptance probability is at most
\(1/2\leq 3/4\).
\end{proof}

\begin{theorem}
\label{QESAT_poly_LPCP}
\lean{QESAT_poly_LPCP}
\uses{QESAT, LPCP, QESAT.verifier}
\leanok
\[
  \mathrm{QESAT}(\F_2,n) \in
  \mathrm{LPCP}[
    \varepsilon_c = 0,
    \varepsilon_s = 3/4,
    \Sigma = \F_2,
    \ell = n + n^2,
    q = 4,
    r = |x| + 2n
  ].
\]
\end{theorem}
\begin{proof}\uses{QESAT.verifier_completeness,QESAT.verifier_soundness}\leanok
% Lecture 08 gives the exact LINEQ randomness as \(m\log|\F|\). In Lean,
% \(|x|\) is the encoded input size and bounds the number of equations \(m\);
% after specializing to \(\F_2\), this yields the displayed \(|x|+2n\) bound.
Use the verifier \ref{QESAT.verifier}. It makes one query for the LINEQ check
and three queries for the tensor self-check, for a total of four. Its
randomness is the LINEQ randomness plus two vectors in \(\F_2^n\). Completeness
and soundness are \ref{QESAT.verifier_completeness} and
\ref{QESAT.verifier_soundness}.
\end{proof}

\begin{theorem}
\label{QESAT_exp_PCP_before_repetition}
\lean{QESAT_exp_PCP_before_repetition}
\uses{QESAT,PCP}
\leanok
\[
  \mathrm{QESAT}(\F_2,n) \in
  \mathrm{PCP}[
    \varepsilon_c = 0,
    \varepsilon_s = 7/8,
    \Sigma = \F_2,
    \ell = 2^{|x|},
    q = O(1),
    r = |x|^{O(1)}
  ].
\]
\end{theorem}
\begin{proof}\uses{QESAT_poly_LPCP,LPCP_to_PCP_ZMod2}\leanok
% The lecture's final no-abstraction slide shares local-correction work in a
% direct QESAT verifier and therefore displays different constants. This
% formalization deliberately routes through the modular LPCP-to-PCP theorem,
% so only the asymptotic constant-query statement is retained here.
Apply \ref{LPCP_to_PCP_ZMod2} to \ref{QESAT_poly_LPCP}. Since the LPCP
soundness is \(3/4\), the converted soundness is
\[
  \max\{7/8,\,3/4+1/100\}=7/8.
\]
The query bound remains constant because \(q=4\), and the randomness is
polynomial in the input size. The Lean formalization pads the proof oracle from
length \(2^{n+n^2}\) to length \(2^{|x|}\).
\end{proof}

\begin{lemma}
\label{QESAT.constant_repetition}
\uses{PCP}
\[\PCP[0,7/8,\F_2,2^{|x|},O(1),|x|^{O(1)}] \subseteq \PCP[0,1/2,\F_2,2^{|x|},O(1),|x|^{O(1)}].\]
\end{lemma}
\begin{proof}
% This repetition is folded into the Lean theorem `QESAT_exp_PCP` rather than
% exposed as a standalone declaration.
Run the \(7/8\)-sound verifier independently six times and accept only if all
six runs accept. Completeness error remains \(0\). For \(x\notin L\), the
acceptance probability is at most \((7/8)^6\leq 1/2\). Query and randomness
bounds are multiplied by the constant \(6\).
\end{proof}

\begin{theorem}
\uses{QESAT,PCP}
\label{QESAT_exp_PCP}
\lean{QESAT_exp_PCP}
\leanok
\[
  \mathrm{QESAT}(\F_2,n) \in
  \mathrm{PCP}[
    \varepsilon_c = 0,
    \varepsilon_s = 1/2,
    \Sigma = \F_2,
    \ell = 2^{|x|},
    q = O(1),
    r = |x|^{O(1)}
  ].
\]
\end{theorem}
\begin{proof}\uses{QESAT_exp_PCP_before_repetition,QESAT.constant_repetition}\leanok
Combine \ref{QESAT_exp_PCP_before_repetition} with the constant repetition
step \ref{QESAT.constant_repetition}.
\end{proof}
